% Reviewer-driven method/results patch for SIGGRAPH Asia draft.
% Date: 2026-05-09
% Scope: problem definition, claim map, method organization, experiment/baseline/metric
% organization, and result placement.
%
% 中文合并说明：
% 1. 本文件不是独立论文，而是可合并到 paper_siga/main.tex 的 patch 草稿。
% 2. 建议优先合并 Problem Definition and Claim Map、Experiment Organization 两部分；
%    Method 中与 main.tex 已重叠的公式可以作为替换或精炼材料。
% 3. 所有关于 texture/PBR 的句子都刻意写成“asset-readiness/export compatibility”，
%    不把纹理渲染当作 topology proof。
% 4. DLA、coral、crystal 统一使用 inspired/scaffold/lattice-like 命名，避免物理模拟 claim。
% 5. 负例保留为 boundary/diagnostic，不作为主结果叙事中心。

\section{Problem Definition and Claim Map}

% 中文注释：这一节建议放在 Method 开头，或者作为 Introduction 最后一段之后的
% formal problem setting。它的目的不是再堆一个大系统 tuple，而是把输入、输出、
% 可以 defend 的 claim，以及不 claim 的东西一次钉死。

We study finite-depth recursive asset synthesis with a frozen native 3D
generator. The input is a root asset $x_0$, an optional appearance or semantic
condition $y$, a recursive program $\mathcal P$, and a recursion budget $D$.
The output is a sequence of projected assets and states
\begin{equation}
  (x_0,z_0), (x_1^\star,z_1),\ldots,(x_D^\star,z_D),
  \qquad
  z_{d+1}\sim
  \mathsf T_{\mathcal P,\theta,d}(\cdot\mid z_d,y,\xi_d),
\end{equation}
where $\theta$ is fixed, $\xi_d$ denotes stochasticity from rule choice,
frontier sampling, or masked generative sampling, and $x_D^\star$ must be
renderable as a mesh or textured mesh. The problem is not one-shot text/image
to 3D generation. It is execution of a recursive program whose intermediate
states must remain attached, addressable, and reusable.

The central design principle is ownership separation. The grammar owns the
recursive structure: handles, transforms, frontiers, attachment certificates,
competition, projection policy, and optional caches. The frozen generator owns
only the native sparse representation, masked local naturalization, decoding,
re-encoding, and selected texture/PBR export. In particular, the generator is
not allowed to globally resample the object topology after each rule.

\begin{table*}[t]
  \centering
  \caption{Claim map used to organize the revised paper. The right column
  states what should be explicitly excluded or moved to supplementary material.}
  \label{tab:reviewer-claim-map}
  \begin{tabular}{@{}p{0.19\textwidth}p{0.31\textwidth}p{0.25\textwidth}p{0.17\textwidth}@{}}
    \toprule
    Claim & Evidence to show in main text & Required metric/protocol & Must not claim \\
    \midrule
    PS-RSLG is a recursive language, not a pipeline &
    compact state, handle-to-proposal rules, projected transition &
    formal semantics and algorithm boxes &
    arbitrary universal procedural language \\
    Per-depth projection stabilizes finite recursion &
    direct vs final-only vs per-depth ablation on conservative programs &
    connected components, LCR, root reachability, orphan mass &
    all operators become topology-clean \\
    Classical systems are covered as restricted domains &
    IFS, L-system, space-colonization, DLA/frontier, shape grammar reductions &
    equations and one comparison table &
    visual dominance over mature procedural systems \\
    Non-tree families are possible when support is grammar-native connected &
    coral/DLA-inspired and crystal/lattice-like connected scaffold results &
    occupancy or surface-voxel connectivity, fixed-camera mesh renders &
    physical DLA, physical crystallization, watertight topology \\
    Texture/PBR export is compatible with selected projected meshes &
    Trellis2 textured GLB renders for selected vine, coral, crystal, baseline assets &
    GLB import/render, PBR availability, surface-voxel diagnostic &
    texture as connectivity proof \\
    Effective recursive refinement is a method behavior &
    depth and parameter rows with zoom crops &
    same guide, same camera, same renderer, changed depth/parameter only &
    calibrated monotonic control unless measured \\
    \bottomrule
  \end{tabular}
\end{table*}

\paragraph{Negative definition.}
The paper should explicitly state what the method does not solve. It does not
guarantee watertight meshes, physical growth simulation, globally consistent
material recursion, or infinite mesh output. Its defensible target is narrower:
training-free finite-depth recursive programs whose sparse support is projected
to an admissible connected state before subsequent rules can fire, with
selected projected meshes exported through a frozen Trellis2 texture/PBR path.

\section{Method Organization Patch}

% 中文注释：本节可用于替换或收紧 main.tex 的 Method 主体。它比旧版大 tuple 更复杂，
% 但复杂性分布在规则、投影、采样、baseline coverage 中，而不是堆在一个状态变量列表里。

\subsection{State, Handles, and Admissible Domain}

At depth $d$, the recursive state is the sparse-latent object
\begin{equation}
  z_d=(\mathcal V_d,\mathbf F_d,\mathcal A_d),
\end{equation}
where $\mathcal V_d\subset h\mathbb Z^3$ is active sparse support,
$\mathbf F_d:\mathcal V_d\rightarrow\mathbb R^q$ is the generator-native feature
field, and $\mathcal A_d$ is a grammar-readable anchor graph. A handle is
\begin{equation}
  h_i=(\sigma_i,T_i,\Omega_i,\alpha_i,\mu_i),
\end{equation}
with type $\sigma_i$, local frame $T_i$, owned region $\Omega_i$, attributes
$\alpha_i$, and edit/material mask $\mu_i$. Auxiliary objects such as traces,
random seeds, cached descriptors, and export metadata are executor state, not
the primary mathematical state.

Connectivity is defined on a sparse neighborhood graph
$G_\eta(\mathcal V_d)$ with root anchors $R_d\subseteq\mathcal V_d$. The attached
support is
\begin{equation}
  \mathcal V_{\mathrm{att}}(z_d)=
  \{v\in\mathcal V_d:\; v\leadsto R_d\ \mathrm{in}\ G_\eta(\mathcal V_d)\}.
\end{equation}
The connectivity violation is
\begin{equation}
  V_{\mathrm{conn}}(z_d)=
  1-\frac{\operatorname{mass}(\mathcal V_{\mathrm{att}}(z_d))}
  {\operatorname{mass}(\mathcal V_d)+\epsilon}.
\end{equation}
We call a state admissible when
\begin{equation}
\begin{split}
  z\in\mathcal Z_{\mathrm{adm}}(\lambda)
  \quad\Longleftrightarrow\quad
  &V_{\mathrm{conn}}(z)\le\epsilon_c,\quad
  V_{\mathrm{frontier}}(z)=0,\quad
  V_{\mathrm{occ}}(z)\le\epsilon_o,\\
  &V_{\mathrm{mat}}(z)\le\epsilon_m,\quad
  |\mathcal V(z)|\le T_{\max}.
\end{split}
\end{equation}
$V_{\mathrm{frontier}}(z)=0$ requires every active growth handle to have a valid
path to the root-attached scaffold. This is the formal reason fragments cannot
be treated as harmless render artifacts: if an orphan component carries an
active handle, it can become a parent, frontier, motif, or material source at
the next depth.

\subsection{Rule Templates and Proposal Semantics}

The recursive language consists of typed handles, local rule templates, and
realization operators. A rule has the form
\begin{equation}
  \rho:\ h_i \mapsto
  \{(\sigma_j,\tau_j,\Omega_j,m_j,\pi_j,\kappa_j,\eta_j)\}_{j=1}^{k}.
  \label{eq:reviewer-rule-template}
\end{equation}
Each emitted item specifies an output type $\sigma_j$, local transform
$\tau_j$, target region $\Omega_j$, edit mask $m_j$, proposal kernel $\pi_j$,
feasibility predicate $\kappa_j$, and optional sampler/material schedule
$\eta_j$. The proposal kernel may grow an attached frontier, split a branch,
attach a sampled exposed cell, transform-copy a motif, refine a terminal
region, or transport a material handle.

The proposal distribution is conditioned on the current sparse state and rule
trace:
\begin{equation}
  P_d\sim
  \prod_{h_i\in\mathcal H_d}
  \prod_{\rho\in\mathcal R(h_i)}
  \pi_\rho(\Delta\mathcal V,\Delta\mathbf F,\Delta\mathcal A
  \mid z_d,h_i,y,\xi_d).
\end{equation}
Accepted proposals must satisfy either direct projectability or a bounded
bridge certificate:
\begin{equation}
  \Delta\mathcal V\subset\mathcal R_{\mathrm{proj}}(z_d;b_d,r_d)
  \quad\mathrm{or}\quad
  \exists B:\ 
  \Delta\mathcal V\cup B\in\mathcal Z_{\mathrm{adm}}(\lambda_d),
  \ \operatorname{cost}(B)\le b_d.
  \label{eq:reviewer-bridge-certificate}
\end{equation}
This condition separates recursive sparse-latent programs from unconstrained
latent editing. A proposed patch is not executable merely because it looks
locally plausible; it must remain connected enough to be used by later rules.

Sparse occupancy competition is a rule-level selection mechanism. For action
$a$ at support location $v$, use
\begin{equation}
\begin{split}
  s(a,v;z_d)=&
  \lambda_{\mathrm{att}}\phi_{\mathrm{att}}(a,v)
  -\lambda_{\mathrm{occ}}\phi_{\mathrm{occ}}(a,v)
  -\lambda_{\mathrm{col}}\phi_{\mathrm{col}}(a,v)\\
  &+\lambda_{\mathrm{front}}\phi_{\mathrm{front}}(a,v)
  +\lambda_{\mathrm{sym}}\phi_{\mathrm{orbit}}(a,v)
  +\lambda_{\theta}\log p_\theta(a,v\mid z_d,y).
\end{split}
\label{eq:reviewer-competition-score}
\end{equation}
The learned term is optional and should be weak in the current experiments. The
main point is that space competition is native to sparse support: occupied
tokens, collisions, attachment paths, and group orbits can be tested before
decoding.

\subsection{Recursive Transition and Masked Generator Use}

The full recursive transition is
\begin{equation}
  z_{d+1}=
  \operatorname{Enc}_{\theta}
  \left[
  \Pi_{\lambda_d}
  \left(
  \operatorname{Dec}_{\theta}
  \left(
  \mathcal T_\theta(z_d;\mathcal R_d,y,\xi_d)
  \right)
  \right)
  \right].
  \label{eq:reviewer-compact-transition}
\end{equation}
The local transition $\mathcal T_\theta$ expands into
\begin{align}
  \mathcal R_d&=\operatorname{Select}(\mathcal A_d,\mathcal R,z_d,\xi_d),\\
  P_d&=\operatorname{Prop}_{\mathcal R_d}(z_d),\\
  \tilde z_{d+1}&=\mathcal M(z_d,\operatorname{Filter}_{\lambda_d}(P_d)),\\
  \bar z_{d+1}&=\mathcal N_\theta(\tilde z_{d+1};m_d,y,\eta_d),
\end{align}
where $\mathcal M$ is masked feature/material merge and $\mathcal N_\theta$ is
the frozen local sampler. This form keeps the main equation short while making
the control flow explicit.

For a flow sampler,
\begin{equation}
  \frac{dz_t}{dt}=v_\theta(z_t,t,y),
\end{equation}
the masked clamp is
\begin{equation}
  z_{t-\Delta t}=
  (1-m_d)\odot z_{\mathrm{anchor}}
  +m_d\odot\Phi_{\theta,t}(z_t,y).
  \label{eq:reviewer-masked-flow}
\end{equation}
Equivalently, for edited sparse features,
\begin{equation}
  \mathbf f_i'=
  (1-\alpha_i)\mathbf f_i^{\mathrm{rule}}
  +\alpha_i\mathbf f_i^{\theta},\qquad i\in m_d.
\end{equation}
The sampler is therefore a local naturalization prior for shape or material. It
is not a global repair engine and it does not decide recursive topology.

\subsection{Projection Operator}

Projection should be described as an explicit algorithm, not only as an ideal
optimization. Given a decoded candidate $x$ and anchors $\mathcal A_d$, build a
mesh, occupancy, or surface-voxel graph $G(x)$. Find root-reachable components,
score orphan components by mass, distance-to-root, bridge cost, collision,
renderability, material consistency, and rule ownership, and then prune,
defer, or bridge them under the schedule $\lambda_d$. Surviving handles are
reattached to projected support; orphan handles and their material descriptors
are removed or marked inactive before re-encoding.

The idealized objective is
\begin{equation}
\begin{split}
  \Pi_{\lambda}(x,\mathcal A)=
  \operatorname*{arg\,min}_{(y,\mathcal A')\in\mathcal X_{\mathrm{adm}}(\lambda)}
  &d_{\mathcal X}(x,y)
  +\eta_{\mathrm{obs}}d_{\mathrm{obs}}(x,y)
  +\eta_m d_{\mathrm{mat}}(x,y)\\
  &+\lambda_{\mathrm{del}}\operatorname{mass}(R)
  +\lambda_{\mathrm{br}}\sum_{\gamma\subset B}\operatorname{len}(\gamma)
  +\lambda_{\mathrm{front}}V_{\mathrm{frontier}}(y,\mathcal A'),
\end{split}
\label{eq:reviewer-projection-objective}
\end{equation}
where $R$ is removed orphan support and $B$ is inserted or certified bridge
support. Prune-only projection is the limiting case
$\lambda_{\mathrm{br}}\rightarrow\infty$. Bridge repair is therefore a
controlled option, not a claim that arbitrary fragmented outputs can be fixed.

\paragraph{Stability statement.}
Let $e(z)$ combine connectivity, orphan-frontier, renderability, and material
violations. If projection suppresses decoded candidates to
$e(\Pi_{\lambda_d}(x))\le\epsilon_P$ and the decode/project/re-encode round trip
adds at most $\epsilon_E$, then the executed states satisfy
\begin{equation}
  e(z_d)\le\epsilon_P+\epsilon_E,\qquad d=1,\ldots,D.
\end{equation}
Final-only cleanup cannot provide this invariant because invalid components may
already have influenced rule selection, frontier sampling, cache updates, and
material propagation before the terminal mesh is cleaned.

\subsection{Classical Procedural Systems as Restricted Domains}

% 中文注释：这一段回答“公式必须 cover IFS/L-system/space colonization/DLA/shape grammar”
% 的批评。它可以放主文精简版，详细证明放 appendix。核心语气是“restricted domains”，
% 不是说当前实现视觉上全面超过传统系统。

The rule template in Eq.~\eqref{eq:reviewer-rule-template} covers standard
procedural families as restricted domains obtained by disabling learned
naturalization, texture export, or projection.

\paragraph{IFS.}
With one patch handle, transform-copy rules, identity sampler, and union merge,
\begin{equation}
  \mathcal V_{d+1}=\bigcup_{i=1}^{k}\tau_i(\mathcal V_d),
\end{equation}
which is Hutchinson iteration on sparse support for contractive transforms.

\paragraph{L-systems.}
Let $\mathcal A_d$ store a word or turtle-frame graph $w_d\in\Sigma^\ast$ and
apply rules synchronously:
\begin{equation}
  w_{d+1}=\prod_{a\in w_d}\rho(a).
\end{equation}
The interpretation operator turns emitted symbols into handles with frame
updates $T_{i+1}=T_i\circ\tau(a_i)$. This recovers deterministic or stochastic
L-systems when $\mathcal N_\theta$ is identity and projection only validates
attachment.

\paragraph{Space colonization.}
With tips $p_i$ and attractors $a\in A$, assign
\begin{equation}
  A_i=\{a\in A:\ i=\operatorname*{arg\,min}_j\|a-p_j\|,
  \ \|a-p_i\|\le r\},
\end{equation}
and propose
\begin{equation}
  p_i'=p_i+\delta
  \frac{\sum_{a\in A_i}(a-p_i)/\|a-p_i\|}
  {\left\|\sum_{a\in A_i}(a-p_i)/\|a-p_i\|\right\|+\epsilon}.
\end{equation}
This is a grow rule with attractor-directed proposal, occupancy exclusion, and
root-attached feasibility.

\paragraph{DLA/frontier growth.}
Let $H(v)$ be a random-walk or learned hitting score on exposed frontier cells.
The attach proposal is
\begin{equation}
  \Pr[\Delta\mathcal V=\{v\}\mid O_d]
  =
  \frac{H(v)\mathbf 1[v\in\partial O_d]}
  {\sum_{u\in\partial O_d}H(u)}.
\end{equation}
The cell is accepted only if adjacent to the occupied root-attached support or
if Eq.~\eqref{eq:reviewer-bridge-certificate} holds. Current experiments should
call this DLA-inspired or frontier-growth scaffolding, not faithful physical DLA.

\paragraph{Shape grammar and symmetry.}
CGA-style split/extrude/repeat rules are contextual handle rewrites over faces,
tiles, or portal frames. A crystal/lattice-like program is represented by a
group or lattice action $G$ over handles and support. If
\begin{equation}
  \forall g\in G,\ \forall\rho\in\mathcal R,\qquad
  g\circ\rho\circ g^{-1}\in\mathcal R,
\end{equation}
and the scheduler expands complete orbits, then the proposed support is
$G$-closed before decoding. Exact equivariance would also require merge,
sampler, projection, codec, and grid quantization to commute with $G$, so the
paper should report symmetry as an approximate metric rather than as a theorem.

\subsection{Material Handles and Texture/PBR Boundary}

Material handles are metadata attached to anchors and masks. A material handle
stores a local intent label, guide condition, PBR-channel availability, and
ownership region. Transform-copy rules transport material handles with the same
local frame as geometry. Growth rules inherit or blend material handles from
parents. Projection deletes material handles whose geometry is pruned and marks
uncertain seams for export validation.

For a projected mesh $x_d^\star$, Trellis2 texture export is
\begin{equation}
  g_d=\operatorname{Tex}_{\theta}(x_d^\star,y_m),
\end{equation}
where $g_d$ is a GLB/PBR asset and $y_m$ is an optional guide. This is an
asset-readiness check for selected projected meshes. It is not structural
evidence. A textured GLB can render well while raw face components, UV islands,
or surface seams remain problematic; therefore topology and material must be
reported with separate diagnostics.

\section{Experiment, Baseline, and Metric Organization Patch}

% 中文注释：建议把现有 Experiments/Results 组织成 claim-first，而不是任务日志式。
% 下面的小节标题可直接成为主文结构，或者作为 main.tex 中 Results 重排的 checklist。

\subsection{Protocol Separation}

Every figure and table should identify which evidence type it supports:
\begin{enumerate}
  \item \textbf{Structural recursion:} occupancy or surface-voxel connectedness,
  root reachability, orphan mass, path-to-root rate, and branch/frontier proxies.
  \item \textbf{Mesh diagnostics:} raw face components, non-manifold edges,
  hole/weld/remesh warnings, Blender import, and renderability.
  \item \textbf{Texture/PBR export:} GLB size, import/render success, material
  guide, PBR-channel availability, seam warnings, and visual QA.
  \item \textbf{Method behavior:} same-condition depth or parameter rows with
  fixed guide, camera, renderer, and texture schedule.
\end{enumerate}
This protocol directly addresses the reviewer concern that texture/PBR quality
cannot substitute for topology evidence.

\subsection{Claim-Based Results Structure}

The revised results section should be ordered as follows.

\paragraph{R1: Projection ablation.}
Compare direct sparse recursion, final-only projection, and per-depth projection
under matched root, operator, depth, renderer, and metric protocol. This is the
main causal evidence. The conservative vine/tree competition rows belong in the
main text; aggressive fork variants should be described as the
stability-expression boundary.

\paragraph{R2: Baseline fairness.}
Treat L-system and space-colonization baselines as strong structural controls,
not strawmen. Compare procedural meshes, direct sparse grammar, final-only
cleanup, and proposed per-depth projection. If texture export is applied to
traditional baselines, present it as a separate diagnostic: successful texture
transfer alone does not prove connected recursive asset readiness.

\paragraph{R3: Non-tree connected scaffold.}
Use crystal/lattice-like and coral/DLA-inspired examples only when mesh renders
and occupancy/surface-voxel connectivity agree. The main text should prefer
grammar-native connected scaffold positives over post-hoc hard-DLA bridge
repair. Hard-DLA bridge experiments are useful boundary evidence if they show
that face-level repair can overstate true support connectivity.

\paragraph{R4: Texture/PBR export.}
Report selected Trellis2 textured GLB exports after projection, with public
guide, schedule, GLB size, render success, and surface-voxel diagnostic. State
that this supports export compatibility and asset inspection, not watertight
topology or material recursion consistency.

\paragraph{R5: Depth and parameter method behavior.}
Show same-condition depth rows and parameter rows for selected cases. These are
not ablations unless only one design choice is changed and all other variables
are matched. Captions should use ``method-behavior visualization'' rather than
``ablation''.

\paragraph{R6: Boundary cases.}
Move fragmented radial, hard-DLA bridge, cache-selected blocky results, and
physically ambiguous cases to a diagnostic or supplementary section unless a
fresh mesh/textured-mesh QA confirms connected visual quality.

\subsection{Metric Definitions to Make Reproducible}

The paper should define the following primary and secondary metrics.

\begin{itemize}
  \item Occupancy connected components: number of components in a voxelized
  occupancy graph under 6-neighborhood connectivity.
  \item Largest component ratio: mass of the largest root-attached component
  divided by total occupied or sampled-surface mass.
  \item Root reachability: fraction of sampled support, branch endpoints, or
  active handles with a path to a root anchor.
  \item Orphan mass: occupied or sampled-surface mass not connected to the root
  component.
  \item Projection survival: fraction of rule-proposed support surviving
  projection at each depth.
  \item Bridge cost: inserted or certified bridge voxels, path length, and
  collision/renderability penalty.
  \item Raw mesh diagnostics: face components, non-manifold edges, holes,
  remesh/weld status, and Blender import/render success.
  \item Texture/PBR diagnostics: GLB size, material guide, PBR voxel/token
  counts when available, import/render success, and seam warnings.
  \item Symmetry/lattice diagnostics: orbit coverage and approximate symmetry
  error, reported as screening metrics rather than theorems.
\end{itemize}

The main structural metric should be occupancy or surface-voxel connectivity.
Raw GLB face components are still important, but they should be labeled as
export diagnostics because material seams and UV islands can fragment the face
graph even when the visible support is connected.

\begin{table*}[t]
  \centering
  \caption{Recommended placement of current result families. The goal is to
  keep the main text claim-aligned and move visually weak or physically
  over-interpretable cases to supplementary diagnostics.}
  \label{tab:reviewer-result-placement}
  \begin{tabular}{@{}p{0.25\textwidth}p{0.28\textwidth}p{0.39\textwidth}@{}}
    \toprule
    Placement & Result families & Rationale and required wording \\
    \midrule
    Main text &
    projection ablation; connected vine/root depth; traditional texture baseline diagnostic; best coral/DLA-inspired connected scaffold; best crystal/lattice-like connected scaffold &
    Directly supports projection, baseline fairness, non-tree connected support, and selected texture/PBR export. Captions must separate structural metrics from texture/export diagnostics. \\
    Supplement or appendix &
    full result matrix; guide sweeps; bismuth cross-guide variants; coral density endpoints; symmetry screening; cache/LOD visible-window demos; extra depth/parameter galleries &
    Useful breadth and method-behavior evidence, but too broad or too weak for the main causal chain unless distilled into one claim-specific figure. \\
    Negative/boundary diagnostics &
    hard-DLA post-hoc bridge; fragmented radial/echo/fork variants; cache-selected blocky DLA/radial/scifi; physically ambiguous crystal/DLA cases &
    Use to justify attachment feasibility, projection placement, and grammar-native connected support. Do not present as final asset success. \\
    Do not use as evidence &
    matplotlib point-cloud previews; uninspected GLBs; partial transfers; texture-only screenshots without mesh/metric QA &
    User-facing and paper-facing qualitative results must be mesh or textured-mesh renders with accompanying diagnostics. \\
    \bottomrule
  \end{tabular}
\end{table*}

\section{Text-Level Replacement Guidance}

% 中文注释：下面这些句子可以直接替换 main.tex 中容易显得 draft/status report 的表达。

\paragraph{Safe topology wording.}
Our strongest structural claim is not that every exported mesh is watertight,
but that connected support must be enforced inside the sparse-latent recursive
transition. Final-only cleanup can improve the terminal mesh, but it cannot
prevent intermediate fragments from becoming active handles, frontiers, cache
entries, or material sources.

\paragraph{Safe baseline wording.}
Classical procedural baselines are strong structural controls under favorable
tube or voxel protocols. We therefore use them as fairness baselines rather
than strawmen, and evaluate our design claim through direct, final-only, and
per-depth projection variants of recursive sparse-latent execution.

\paragraph{Safe non-tree wording.}
Crystal and coral examples are reported as connected scaffold families. They
demonstrate group/lattice-style and coral/DLA-inspired recursive supports with
selected PBR export, not physical crystallization, faithful DLA simulation, or
watertight topology.

\paragraph{Safe texture wording.}
Textured GLB export is evaluated separately from structural connectivity.
Surface-sampled voxel connectivity and fixed-camera mesh renders support
connected-scaffold claims, while raw face components and material seams remain
export diagnostics.

\paragraph{Words to remove from main paper.}
Avoid ``draft'', ``candidate'', ``smoke test'', ``technically compatible'',
``current paper use'', ``bug fixed'', and ``failed export'' in the main text.
Those phrases belong in internal notes or supplementary diagnostics, not in the
submission narrative.
